% Sección 4
\section{Las preguntas de la mentoría}

La mentoría dejó las cinco preguntas que siguen. Dos agentes investigan esta noche lo que falta para responderlas. Uno busca alternativas, actores y el marco oficial de recuperación. El otro reúne lo que dijeron los actores en redes, prensa e informes de gestión. Lo que depende de ellos lleva la marca \eninvestigacion.

\begin{enumerate}[leftmargin=1.6em, itemsep=0.1em]
\item ¿Por qué hacemos esto y para qué sirve?
\item ¿A quién le sirve y a quién impacta?
\item ¿Hay algo parecido? ¿Por qué no se usa en Manizales o en la región?
\item ¿Qué prototipo va a medir el efecto en la región?
\item ¿Cómo es el piloto de 30 a 90 días para la recuperación del territorio?
\end{enumerate}

\subsection{Por qué hacemos esto y para qué sirve}

\paragraph{Lo que sabemos} La comida llegó a Manizales, y la dificultad estuvo en repartirla. ABACO movió más de 2.000 toneladas en la primera semana y nombró la llegada al territorio como su reto principal. En Bogotá, el acopio de El Campín reunió carga más rápido de lo que podía despachar. Los cortes de energía dejaron por fuera a los perecederos, y dos vías de salida seguían cerradas dos días después del sismo.

DespachoX sirve para responder cada mañana la pregunta de quien despacha. Con lo que hay en bodega y los vehículos disponibles, esa persona tiene que decidir qué le manda a quién, en qué orden y por dónde. El resultado que buscamos es que, con la misma comida y los mismos vehículos, más pedidos urgentes lleguen a tiempo y menos alimentos se dañen. Esos son los dos indicadores del reto.

La trazabilidad responde a otra necesidad, la confianza. En Colombia hay antecedentes de ayudas que se perdieron o se desviaron, como los 80 carrotanques que la UNGRD compró para La Guajira y que terminaron deteriorados sin uso. En 2026 también hubo preguntas públicas sobre el destino de las donaciones. Un código por paquete, confirmado al entregarse, permite mostrar a quien donó que su ayuda llegó.

\paragraph{Lo que sigue abierto} La ficha del reto dice que el desabastecimiento llega «en cuestión de 48 horas», y no encontramos fuente que lo respalde para este sismo. Por eso no lo usamos como hecho. Tampoco sabemos si en Manizales se vencieron alimentos donados, ni cuántas familias recibieron ayuda alimentaria en la ciudad. La pregunta de fondo, cómo despacha hoy el Banco de Alimentos de Manizales, está \pendiente.

\subsection{A quién le sirve y a quién impacta}

\paragraph{Lo que sabemos} La ficha del reto nombra como usuarios a la Cruz Roja, la Defensa Civil, el Banco de Alimentos, los transportadores y los comités barriales. El equipo pone en el centro a la persona que despacha. Puede ser una funcionaria del Banco de Alimentos con años de experiencia o un voluntario que llegó ayer, y en ambos casos decide con información incompleta. Le siguen el líder de barrio, que pide en nombre de sus vecinos, y el conductor, que conoce el estado real de las vías mejor que cualquier mapa. Las familias no van a usar la herramienta. Lo que notan es si la comida llegó a tiempo y en buen estado.

La evidencia señala a tres grupos de afectados en particular. Están las personas en albergues como el Coliseo Mayor y las familias autoalbergadas, que ABACO puso entre sus prioridades. Están también las familias rurales dispersas que quedaron fuera del Registro Único de Damnificados.

El sistema tiene un costo que ya medimos en simulación. Al atender primero lo crítico, el punto peor atendido espera más que con el despacho por orden de llegada. Por eso el sistema propone y una persona aprueba.

\paragraph{Lo que sigue abierto} Qué actores operaron de verdad en Manizales y qué dijeron está \eninvestigacion. Qué comunas concentraron a los damnificados sigue \pendiente. Si el Banco de Alimentos de Manizales y la Cruz Roja de Caldas se reconocen como usuarios también está \pendiente.

\subsection{Si hay algo parecido y por qué no se usa aquí}

\paragraph{Lo que sabemos} Existen piezas sueltas. Hay herramientas abiertas que calculan rutas con vías bloqueadas, como el openrouteservice para desastres de HeiGIT, VROOM y Valhalla. Otras registran entregas, como LMMS de World Vision, SCOPE del Programa Mundial de Alimentos (PMA) y KoBo, y otras recogen reportes ciudadanos, como Ushahidi. La OPS creó en los años noventa un sistema regional de manejo de suministros humanitarios, SUMA/LSS, que Colombia incorporó en su normativa.

Después del sismo aparecieron plataformas para informar y para recaudar. Cuidar a Colombia, de la Universidad EAFIT, publicó un directorio verificado con 58 puntos de acopio en 32 ciudades. Un bot de WhatsApp de JumpCube con AWS reunió información oficial para los damnificados. ABACO recibió donaciones en línea y Rappi habilitó un botón para donar.

La región tiene además antecedentes útiles. En Haití, en 2010, Mission 4636 convertía mensajes de texto en reportes estructurados y ubicados en el mapa en unos cinco minutos. En Brasil, en 2024, la web Salva RS! se construyó en unas 48 horas para ordenar los pedidos de rescate que circulaban por WhatsApp, y la Defensa Civil la usó desde el tercer día. En México, en 2017, Verificado19S comprobaba cada reporte y publicaba necesidades concretas por zona.

Lo que no encontramos es una herramienta abierta que una todo el recorrido. Ese recorrido va del mensaje de un vecino a la ruta de un vehículo, con el cierre que reportó un conductor y el vencimiento de lo que va en la carga. Las plataformas de 2026 que encontramos se dedican a informar o a recaudar, y ninguna de ellas despacha pedidos.

Nuestra explicación de por qué algo así no se usa en Manizales es que no existe unido. Cada pieza resuelve una parte, y juntarlas bajo presión supera lo que un equipo de emergencia puede montar en pocos días. Esta explicación es inferencia nuestra. De SUMA/LSS no sabemos si alguien lo usa en Caldas, porque la investigación no verificó su uso por parte de la comunidad.

\paragraph{Lo que sigue abierto} Qué alternativas usan hoy la Alcaldía, la UGR y los organismos de socorro está \eninvestigacion. Si el Banco de Alimentos de Manizales usa algún sistema propio está \pendiente.

\subsection{Qué prototipo va a medir el efecto en la región}

\paragraph{Lo que sabemos} El prototipo es una mesa de despacho. Los comités barriales piden por chat, con texto o con voz, y los conductores reportan las vías bloqueadas por el mismo canal. Un modelo de lenguaje lee y estructura cada mensaje. Luego un optimizador asigna los pedidos a una flota mixta y calcula las rutas con lo que hay en bodega y lo que vence primero. El despachador aprueba cada plan, y cada paquete lleva un código QR hasta que alguien confirma la entrega. La regla del equipo es que el modelo lee, el optimizador calcula y la persona decide.

El 20 de noviembre el sistema corrió por primera vez de punta a punta en la nube. Un mensaje recibido por el bot se convirtió en un pedido confirmado y en un plan calculado. En esa prueba nadie aprobó el plan ni probó un cierre.

Hoy medimos los dos indicadores del reto solo en simulación, y la sección~\ref{sec:prototipo} da las cifras. En un piloto, la confirmación de cada entrega con el código QR daría los datos para medirlos en la operación real. Esto último es inferencia nuestra.

\paragraph{Lo que sigue abierto} No hay todavía ninguna medición con datos reales de Manizales. La simulación usa supuestos nuestros sobre demanda, flota, inventario y cierres, y no mide el valor de recibir los pedidos por chat ni el de la trazabilidad.

\subsection{Cómo es el piloto de 30 a 90 días}

\paragraph{Lo que sabemos} Tenemos una lista de aliados posibles y lo que se le pediría a cada uno. ABACO y el Banco de Alimentos de Manizales serían los usuarios del piloto y podrían aportar datos de los despachos de agosto. A la UGR de Manizales se le pediría la lista de albergues y de vías cerradas, y el aval para operar. La Cruz Roja de Caldas podría probar el tablero en el reparto. El equipo humanitario de OpenStreetMap (HOT) reporta una ortofoto de dron de Manizales del 19 de agosto, y el IDEA de la UNAL tiene datos de lluvia. Exportar la trazabilidad en un formato compatible con SUMA/LSS sería un buen hito para el día 60.

Del marco oficial conocemos una pieza. El Fondo Milagro, creado por el Decreto Legislativo 1379 del 9 de septiembre, financia la reconstrucción. Lo administra Fiducoldex y recibió su primer billón de pesos el 7 de octubre. Su objeto son la demolición, la vivienda y los colegios, y la ayuda alimentaria queda por fuera.

\paragraph{Lo que sigue abierto} El diseño del piloto se escribe esta noche (sección~\ref{sec:piloto}). El marco oficial de recuperación, más allá del Fondo Milagro, está \eninvestigacion. Ningún aliado ha confirmado que nos recibiría, así que ese punto está \pendiente.
